% Indentation
\usepackage{indentfirst}
\setlength{\parindent}{30pt}

% Section titles formatting
\usepackage{titlesec}
\titleformat{\section}
    {\normalfont\bfseries\fontsize{14}{20}\selectfont}
    {\thesection}{1em}{}
\titlespacing*{\section}
    {0pt}{10pt}{0pt}

\titleformat{\subsection}
    {\normalfont\bfseries\fontsize{13}{20}\selectfont}
    {\thesubsection}{1em}{}
\titlespacing*{\subsection}
    {0pt}{10pt}{0pt}

\titleformat{\subsubsection}
    {\normalfont\bfseries\fontsize{12}{20}\selectfont}
    {\thesubsubsection}{1em}{}
\titlespacing*{\subsubsection}
    {0pt}{10pt}{0pt}

% Headers and Footers
\usepackage{fancyhdr}
\pagestyle{fancy}
\fancyhead{}
\fancyfoot{}
\fancyhead[L]{\textsl{\leftmark}}
\fancyhead[R]{\headerauthor}
\fancyfoot[C]{\thepage}
\renewcommand{\headrulewidth}{0.4pt}
\renewcommand{\footrulewidth}{0pt}

% Define placeholders for metadata used in header/footer
\providecommand{\headerauthor}{}

% Long inline-code identifiers (e.g. long_column_names) have no natural break
% point and can be wider than their table cell, overflowing into the next
% column. \texttt disables hyphenation, so inside longtables only, let code
% spans wrap at any character via seqsplit; prose elsewhere is unaffected.
\usepackage{seqsplit}
\let\originaltexttt\texttt
\newcommand{\seqsplittexttt}[1]{\seqsplit{#1}}
\AtBeginEnvironment{longtable}{\let\texttt\seqsplittexttt}
\AtEndEnvironment{longtable}{\let\texttt\originaltexttt}
